\iflongAllModels
\section{Long Autonomous Rollouts: Four-Model Comparison}
\else
\section{Long Autonomous Rollouts: Completed Maps, Pending Solver References}
\fi
\label{app:long-horizon}

\paragraph{Test case and current scope.}
We extend the selected seed-zero maps to 4,096 successive forecast intervals from one held-out initial history: trajectory 180, frames 0--4. Each model carries eight independent generated histories after initialization. The It\^o map uses eight generative steps per frame and the flow map one. \iflongAllModels Both maps and the matched 64-step stochastic-interpolant and flow-matching references have completed all 4,096 intervals.\else Both maps are complete. The matched 64-step stochastic-interpolant and flow-matching trajectories are not yet complete, so this appendix reports a \emph{partial model comparison}.\fi It is separate from the all-20-trajectory, three-student confirmation in the main text.

\paragraph{Time and reference data.}
The recovered companion archive supplies a mean saved-time increment of 0.49989997016; thus this horizon spans approximately 2,047.59 nominal simulation-time units. The primary coordinate remains stored snapshot intervals. Exact matching of all coarse fields to the companion's stride-two downsampling establishes the timestamp association, but does not identify the original viscosity or forcing parameters. Only 95 future truth fields are available, corresponding to archive frames 5--99. RMSE and spread/skill stop at lead 95; later diagnostics do not measure accuracy against a reference continuation.

\paragraph{Diagnostics.}
We inverse-transform each member into the archive's vorticity units, without changing the internal model trajectory. For $P$ spatial cells, member $e$ has area-mean enstrophy $Z_e=(2P)^{-1}\sum_p\omega_{ep}^2$. We report its ensemble mean, rather than the generally smaller enstrophy of the ensemble-mean field. In fact,
\begin{equation}
\frac1E\sum_e Z_e
= Z(\bar\omega)+\frac{E-1}{2E}S_{\mathrm{RMS}}^2,
\end{equation}
where $S_{\mathrm{RMS}}$ uses unbiased sample variance as in Appendix~\ref{app:metrics}. Corrected spread is $\sqrt{(E+1)/E}\,S_{\mathrm{RMS}}$; the ratio divides it by ensemble-mean RMSE. No amplitude clipping, rescaling, truth reset or replacement of a failed member is allowed. \longAuditSentence

\begin{table*}[t]
\centering
\caption{Completed map trajectories only; solver references remain pending. Late statistics average leads 3,585--4,096. $Q_{1:95}$ is the ratio of pooled corrected RMS spread to pooled RMSE over the available-reference interval; it is not the mean of per-lead ratios. All eight members remain finite through lead 4,096.}
\label{tab:long-maps}
\small
\begin{tabular}{lrrrrr}
\toprule
Model (steps/frame) & Mean $Z$, leads 1--95 & Late mean $Z$ & Late mean $S_{\mathrm{RMS}}$ & $Q_{1:95}$ & Max. $|\omega|$ \\
\midrule
It\^o map (8) & 1.9777 & 1.4047 & 1.6605 & 0.6709 & 13.129 \\
Flow map (1) & 3.8592 & 3.8366 & 2.7679 & 0.9387 & 16.730 \\
\bottomrule
\end{tabular}
\end{table*}


\paragraph{Observed behavior.}
Both maps remain finite through all 4,096 intervals, but their amplitude and uncertainty differ substantially (Table~\ref{tab:long-maps}). The It\^o trajectory loses member enstrophy and approaches a lower-amplitude range. Over the final 512 intervals its mean enstrophy is 1.4047, compared with 3.8366 for the flow map. The contextual mean across 2,000 stored held-out fields is 4.6654; this archive statistic is not an independently verified long-time target distribution. More directly, mean enstrophy of the available paired truth over leads 1--95 is 4.3526, whereas the corresponding It\^o and flow-map values are 1.9777 and 3.8592. Thus the It\^o amplitude loss is already visible where reference data exist.

The ratio of pooled corrected spread to pooled RMSE over those 95 leads is 0.6709 for the It\^o map and 0.9387 for the flow map (Figure~\ref{fig:long-maps-skill}). Their pooled ensemble-mean RMSE values are relatively close, 3.0148 and 3.0789, despite the spread difference. Low error of an ensemble mean therefore does not establish retained forecast variability. These correlated leads from one initial condition do not provide an independent-case calibration estimate. The observation supports numerical boundedness of these two runs, with marked It\^o dissipation; it does not establish accurate invariant statistics or stability for other initial conditions. \longTeacherComparison No mechanistic attribution or new timing-based speedup claim is made from this extension.

\paragraph{Spatial enstrophy spectra.}
To resolve the scales of the first and last generated fields, let
$\widehat\omega_e(q)=P^{-1}\sum_p\omega_{ep}\exp(-2\pi i\,q\cdot p/n)$
for the $n\times n$ grid, and group integer wavevectors by nearest radial shell. The shell-summed enstrophy is
\begin{equation}
\mathcal Z_e(k)=\frac12\sum_{q:\,\lfloor\lVert q\rVert+1/2\rfloor=k}
|\widehat\omega_e(q)|^2.
\end{equation}
Parseval's identity gives $\sum_{k\ge0}\mathcal Z_e(k)=Z_e$.
This is a sum over modes in each shell, not a radially averaged power spectrum; it does not redefine the earlier spectral-error metric. The spectra describe the observed $128^2$ grid without additional filtering. Frequencies are cycles per grid domain; no physical domain length is assumed. Figures~\ref{fig:long-spectrum-first} and~\ref{fig:long-spectrum-last} show separate first- and last-step spectra on log--log axes. The zero-frequency component is retained in the audit but omitted from these logarithmic axes. All square-grid modes are included; shells above the axis Nyquist frequency $k=64$ contain only corner modes. No spectral values are floored. Reference truth is available for the first forecast step only. At lead 4,096 the It\^o map has lower mean shell enstrophy than the flow map in every shell $1\le k\le64$. This resolves their amplitude difference by spatial scale, but does not identify a correct terminal spectrum without a matching reference.

\begin{figure*}[t]
\centering
\includegraphics[width=.82\textwidth]{figures/first_step_enstrophy_spectrum.pdf}
\caption{Spatial enstrophy spectrum at the first forecast step (lead 1). Curves are means of eight member shell-summed spectra, with member 5--95\% bands, in archive vorticity units. The black curve is the matching observed field. Both axes are logarithmic; $k$ denotes radial spatial frequency in cycles per grid domain. The dotted line marks the axis Nyquist frequency 64; higher shells contain square-grid corner modes. \longScopeCaption}
\label{fig:long-spectrum-first}
\end{figure*}

\begin{figure*}[t]
\centering
\includegraphics[width=.82\textwidth]{figures/last_step_enstrophy_spectrum.pdf}
\caption{Spatial enstrophy spectrum at the last forecast step (lead 4,096), using the same normalization and log--log spatial-frequency axes as Figure~\ref{fig:long-spectrum-first}. Member-mean spectra and 5--95\% member bands describe generated fields only: no reference continuation is available at this lead. Differences between maps therefore show different retained spatial variability, not verified long-time spectral accuracy. \longScopeCaption}
\label{fig:long-spectrum-last}
\end{figure*}

\begin{figure*}[t]
\centering
\includegraphics[width=\textwidth]{figures/long_horizon_maps_skill.pdf}
\caption{Available-reference diagnostics for the same trajectories, restricted to leads 1--95. Left: corrected spread (dashed) and ensemble-mean RMSE (solid). Right: their per-lead ratio. The correction is $\sqrt{9/8}$ applied to unbiased RMS spread for eight members. Similar RMSE does not compensate for the It\^o map's smaller spread. No skill values are extrapolated beyond the stored reference. \longScopeCaption}
\label{fig:long-maps-skill}
\end{figure*}
